\section{基于xv6的RL-scheduler源码实现}
接下来解少Qleanring再系统内核中的具体实现
\subsection{总体框架设计}
xv6的主要逻辑主要集中在procscheduler中，其决策过程在原有原有scheduler的基础上增加了，全局状态构建，状态选择，根据当前状态选择action，并且通过action策略选择下一个满足条件的进程，然后进行上下文切换运行该进程。
\todo{一个说明过程的流程图}
\todo{一个item来列举整个流程}
\todo{具体代码来说明}
更加具体来说，整个RL-scheduler主要分为如下三个部分：数据结构、决策机制和学习机制。

\subsection{数据结构}
从系统角度来看，数据结构的设计主要是为整个RL-scheduler提供系统状态的相关信息的统计与维护。
其中因为xv6操作系统原本设计是作为一个简单的教学操作系统，所以其对浮点数的计算的指令不支持，在实际设计的时候通过定义SCALE常数来对整数进行放缩来支持有理数运算。
\todo{在原来Qlearning的计算公式上，给出加入SCALE以后的计算公式}
\subsection{进程统计字段}
进程作为scheduler调度的基本单位，通过构造进程的状态，使得scheduler能够获取整个调度队列的状态，从而能够更好的进行决策。
更加具体来说，在实现中，在struct proc中设计增加了一些新的统计字段，并且进行了维护。
\todo{进程统计字段列举说明+ 更新说明 }
\subsection{Q表}
Qlearning算法中需要Q表来存储$Q(s,a)$的状态值，由于xv6的内核环境资源有限，本次实验采用了二维数组
int qtable[NSTATE][NACTION]来存储。
\todo{Q表更加具体定义}
\subsection{状态编码映射}
由于xv6的内核环境资源有限，追求轻量化的目标，将proc中的统计字段通过离散化映射为Q表中的状态。
\todo{状态编码的细节补充}
\subsection{决策机制}
从算法的角度来看，决策主要依赖与当前全局状态的编码，根据当前的全局状态，scheduler会采用$\epsilon$-贪心的策略根据Q表，进行调度策略的选择（action），并且使用该调度策略选择下一个满足条件的进程。
其中因为$\epsilon$-贪心策略需要使用到随机数，而xv6本身并不支持，所以添加了随机数模块。
\todo{随机数模块的简单介绍}
\subsection{学习机制}
当整个系统进行完决策，选择进程由于各种原因上下文切换回scheduler中后，此时已经经过了一个隐式的系统状态变换，所以整个调度器中会再进行一次系统状态的编码获取当前系统状态，并且依据前后系统状态变化按照给定的reward的策略，进行reward计算，最终更新Q表，进行下一轮循环。
\todo{reward计算细节}